%% kapa: exam=1 round=35 year=2024 period=1 form=A subject=civil_law subject_ko=민법 seq=11 number=11
%% source: Q-Net 감정평가사 기출문제 (cst00309) artlSeq=5213229 file=raw/1008/5213229/1교시 문제지 원본.pdf
%% conversion: kordoc 4.24.51 -> tools/gen1.py (auto); review=proofread-sonnet; flags=none
\begin{question}{11}
사기ㆍ강박에 의한 의사표시에 관한 설명으로 옳은 것은? (다툼이 있으면 판례에 따름)
\begin{choices}
  \item 피기망자에게 손해를 가할 의사는 사기에 의한 의사표시의 성립요건이다.
  \item 상대방이 불법으로 어떤 해악을 고지하였다면, 표의자가 이로 말미암아 공포심을 느끼지 않았더라도 강박에 의한 의사표시에 해당한다.
  \item 상대방의 대리인이 한 사기는 제3자의 사기에 해당한다.
  \item 단순히 상대방의 피용자에 지나지 않는 사람이 한 강박은 제3자의 강박에 해당하지 않는다.
  \item 매도인을 기망하여 부동산을 매수한 자로부터 그 부동산을 다시 매수한 제3자는 특별한 사정이 없는 한 선의로 추정된다.
\end{choices}
\end{question}
\answer{5}
